\begin{pastexam}[label={px:ccu-114-swa-05}]{中正資工所甲組 114 學年度 第 5 題}
（單選）

(5\%) (Choose \underline{one} best answer) Consider the integer array \{52, 21, 40, 88, 33, 57, 46, 71\}.
Using the two-pointer version of quicksort, what will the array look like after moving the first pivot
(i.e., 52) to its correct position? Note that the first and the second pointers start at the positions of
21 and 71, respectively.
\begin{enumerate}[label=(\alph*)]
\item \{21, 40, 33, 46, 52, 71, 57, 88\}
\item \{21, 33, 40, 46, 52, 57, 71, 88\}
\item \{33, 21, 40, 46, 52, 57, 88, 71\}
\item \{46, 21, 40, 33, 52, 57, 88, 71\}
\item None of the above.
\end{enumerate}
\end{pastexam}

\begin{pastsoln}
答案是 (c)。

雙指標的分割以第一個數 $x = 52$ 為基準：左指標 $i$ 從 $21$ 往右找第一個 $\ge x$ 的數，右指標 $j$ 從 $71$
往左找第一個 $\le x$ 的數；$i < j$ 就交換兩者、兩個指標各走一步，繼續找；$i \ge j$ 時停止，把基準和
$A[j]$ 交換，基準就到了正確位置。 這是資結書第~14.5~節 Hoare 分割的做法，差別只在最後把基準換到中間。 位置從 $1$ 編號。

\begin{center}
\small
\begin{tabular}{l*{8}{c}l}
\toprule
& 1 & 2 & 3 & 4 & 5 & 6 & 7 & 8 & 說明 \\
\midrule
開始 & 52 & 21 & 40 & 88 & 33 & 57 & 46 & 71 & $i = 2$，$j = 8$ \\
找 & & & & $i$ & & & $j$ & & $21, 40 < 52$，$i$ 停在 $88$；$71 > 52$，$j$ 停在 $46$ \\
交換 & 52 & 21 & 40 & 46 & 33 & 57 & 88 & 71 & $i = 5$，$j = 6$ \\
找 & & & & & $j$ & $i$ & & & $33 < 52$，$i$ 停在 $57$；$57 > 52$，$j$ 停在 $33$ \\
歸位 & 33 & 21 & 40 & 46 & 52 & 57 & 88 & 71 & $i \ge j$，交換 $A[1]$ 和 $A[j] = A[5]$ \\
\bottomrule
\end{tabular}
\end{center}
結果 $\{33, 21, 40, 46, 52, 57, 88, 71\}$，$52$ 左邊都比它小、右邊都比它大，就是 (c)。

其他選項：(a)(b) 的左半段已經排好或重排過，一次分割做不到；(b) 是整個排好的結果。 (d) 的 $52$ 也在第 $5$
格、兩邊也分好了，但 $46$、$33$ 的位置不對：第一次交換是 $88$ 和 $46$，$46$ 留在第 $4$ 格，最後和基準交換、
換到第 $1$ 格的是 $33$。 (c) 正確，所以 (e) 不選。
\end{pastsoln}
